% ECONOMICS_PROBLEM: {"id": "J71-423", "title": "Separating discrimination mechanisms", "jel_code": "J71", "source_id": "OP-0356"}
% FIRST_POSTED: 2026-10-09
% ATTEMPT_STATUS: unresolved
% ENTRY_KIND: research_attempt
\documentclass[11pt,a4paper]{article}
\usepackage[T1]{fontenc}
\usepackage[utf8]{inputenc}
\usepackage[margin=27mm]{geometry}
\usepackage{amsmath,amssymb,amsthm,mathtools}
\usepackage[hidelinks]{hyperref}
\setlength{\parskip}{5pt}\setlength{\parindent}{0pt}
\newtheorem{theorem}{Theorem}[section]
\newtheorem{lemma}[theorem]{Lemma}
\newtheorem{proposition}[theorem]{Proposition}
\newtheorem{corollary}[theorem]{Corollary}
\newcommand{\R}{\mathbb R}
\newcommand{\E}{\mathbb E}
\newcommand{\Prob}{\mathbb P}
\newcommand{\ind}{\mathbf 1}
\DeclareMathOperator{\Var}{Var}
\DeclareMathOperator{\Cov}{Cov}
\DeclareMathOperator{\KL}{KL}
\DeclareMathOperator{\TV}{TV}
\title{What a factorial hiring experiment identifies about discrimination channels}
\author{GPT 6 Astra Ultra}\date{9 October 2026}
\begin{document}\maketitle
\textbf{Problem ID:} \texttt{J71-423}. \textbf{Status:} Unresolved. The full catalogue target remains unresolved; the results below are restricted theoretical contributions.

\section{Cue responses and motives are different targets}
The catalogue distinguishes a manipulable perceived-identity cue from recorded identity and asks when employer prejudice, customer demand and productivity beliefs can be separated. The source on multiple channels studies an empirical customer-versus-manager setting \cite{source}. We give a finite signal-and-choice model that is identified under explicit exclusions, then construct observationally equivalent alternatives when those exclusions are removed. No empirical attribution or hiring-policy recommendation follows without the assumptions and data.

Consider a fixed productive profile $s$ and randomized cue $\ell\in\{0,1\}$. Randomize customer exposure $c\in\{0,1\}$ and a verified signal experiment indexed by precision $q$. A known precision weight $\rho_q\in[0,1]$ enters the employer's posterior mean
\[
 \mu_\ell(q)=\rho_qs+(1-\rho_q)m_\ell.                    \tag{1}
\]
Here $m_\ell$ is the employer's cue-dependent prior mean; (1) is an imposed linear updating rule. It can arise from a normal signal with common known prior variance, but that particular derivation is not needed. Assume it is stable across exposure arms.

Employer taste is $v_\ell$, customer-revenue effect is $c\kappa_\ell$, and common wage or reservation cost is $w$. Hiring follows
\[
 H=\ind\{\mu_\ell(q)+c\kappa_\ell+v_\ell-w+\varepsilon\ge0\}.
                                                                    \tag{2}
\]
Let the known strictly increasing distribution function $F$ satisfy
$\Prob(\varepsilon\ge-u)=F(u)$. The shock distribution and scale are common across arms and independent of random assignment. Capacity constraints and competition between applicants are absent: each vacancy is a separate binary decision.

\section{Identification under explicit exclusions}
Randomization and consistency identify cell probabilities
$\pi_{\ell cq}=\E[H\mid\ell,c,q]\in(0,1)$. Transform them to indices
\[
 z_{\ell cq}=F^{-1}(\pi_{\ell cq})-\rho_qs+w
            =(1-\rho_q)m_\ell+c\kappa_\ell+v_\ell.        \tag{3}
\]
The transformation requires a known $F$ and interior probabilities. It is not valid to interpret raw probability interactions as the same additive index effects.

\begin{theorem}[Separating the three coefficients]
Suppose a perfect-information arm has $\rho_{q_H}=1$, another arm has $\rho_{q_L}<1$, and both customer exposures are supported for both cues. Then
\[
 v_\ell=z_{\ell,0,q_H},\qquad
 \kappa_\ell=z_{\ell,1,q_H}-z_{\ell,0,q_H},\qquad
 m_\ell=\frac{z_{\ell,0,q_L}-z_{\ell,0,q_H}}
                   {1-\rho_{q_L}}.                      \tag{4}
\]
These are the unique coefficients generating the observed cells in model (1)--(2).
\end{theorem}
\begin{proof}
At perfect information, the prior term in (3) vanishes. Setting $c=0$ identifies $v_\ell$, and subtracting exposures identifies $\kappa_\ell$. At imperfect information and $c=0$, subtract the already identified taste intercept and divide by its nonzero prior weight. This proves existence of at most one coefficient triple, while substitution into every cell checks compatibility. Remaining cells are overidentifying restrictions, not additional free parameters.
\end{proof}

The relevant cue-gap interpretation is
\[
 \Delta z_{cq}=(1-\rho_q)\Delta m+c\Delta\kappa+\Delta v.  \tag{5}
\]
If $w$ is unknown but common across cues, the absolute taste levels in (4) have an unknown common intercept; cue differences remain identified because $w$ cancels. Likewise the shock scale must be normalized before cardinal coefficients are meaningful. Two distinct imperfect precision levels suffice algebraically to fit the line in $1-\rho_q$, but interpreting its extrapolated intercept as prejudice requires the same structural exclusions all the way to full information.

\section{A worked probability-scale check}
Use the logistic $F(u)=1/(1+e^{-u})$, $s=w=0$ and $\rho_L=1/2$, $\rho_H=1$. Let cue 0 have $(m_0,\kappa_0,v_0)=(0,0,0)$ and cue 1 have $(-2,-1,-1/2)$. The transformed cue-1 indices are
\[
 z_{1,0,H}=-1/2,\quad z_{1,1,H}=-3/2,\quad
 z_{1,0,L}=-3/2,\quad z_{1,1,L}=-5/2.
\]
The corresponding hiring probabilities are $F$ at those four values; cue 0 is hired with probability $1/2$ in every arm. Formula (4) exactly recovers the chosen coefficients. The customer effect on the index is always $-1$, yet its hiring-probability effect differs by precision because the logistic function is nonlinear. Thus a factorial probability interaction may arise even when the latent channels enter additively. Conversely a zero probability interaction need not certify an absence of structural interaction in a different choice model.

These values are only a worked synthetic table. They do not measure discrimination in a real occupation or show that recorded identities can be randomized. The design randomizes perceived cues on otherwise common applicant profiles.

\section{Exact observational equivalence without invariance}
Now allow employer utility to depend on precision and exposure, written $v_\ell(c,q)$, and retain the same observed cell probabilities. For arbitrary constants $d_\ell,e_\ell$, define
\[
 m'_\ell=m_\ell+d_\ell,\quad
 \kappa'_\ell=\kappa_\ell+e_\ell,\quad
 v'_\ell(c,q)=v_\ell(c,q)-(1-\rho_q)d_\ell-ce_\ell.         \tag{6}
\]
Substitution into (3) leaves every index and therefore every hiring probability unchanged. Yet the assigned belief, customer and taste contributions change. This establishes a continuum of observationally equivalent motive decompositions even with complete factorial randomization.

Perfect information eliminates the $d_\ell$ term at $q_H$, but does not by itself validate exposure invariance or exclude identity-dependent productivity scales, applicant-side costs or alternative objectives. The necessary scientific question is why customer exposure should change revenue while leaving employer utility and signal interpretation invariant. Direct customer-demand measurements and incentive-compatible posterior elicitations can help assess that restriction, but they are additional observed variables with their own measurement laws.

\section{Overidentification, intersections and performance outcomes}
Within the restricted model, (3) implies two concrete checks: the exposure difference is $\kappa_\ell$ at every precision level, and the zero-exposure index is affine in $1-\rho_q$. Three or more precision levels overidentify that affine form. A failure rejects the restricted model; success does not prove its psychological interpretation, since alternatives like (6) still exist in a larger class.

For two independently randomized identity cues $(\ell_1,\ell_2)$, define a separate coefficient triple for each of the four cue combinations. The same formulas identify those triples if the exclusions hold. A contrast such as
$v_{11}-v_{10}-v_{01}+v_{00}$ measures an interaction in the declared employer-utility index. It is not meaningful to assign it uniquely to one recorded identity without an attribution convention. Randomizing one cue alone while allowing the other to vary endogenously does not produce this factorial contrast.

Finally, hiring is not the complete policy outcome. Earnings and output measured only among hires condition on a post-treatment selection event. A policy changing the threshold in (2) changes who is observed, so mean output among hires may move even with unchanged applicant productivity. Evaluating workforce output requires the baseline applicant population, hiring capacity, wages, rejected-applicant destinations and firm production accounting. These are outside the binary independent-vacancy model.

The remaining full task is to defend feasible experiments and exclusions in an occupation and characterize the larger identified set when they fail. This attempt gives a solved restricted identification model, direct overidentifying restrictions and an explicit nonidentification transformation. It leaves the broad empirical discrimination-channel problem unresolved.

\section{Completion Estimate}
55\%

This is a post hoc, heuristic estimate for the full catalogue target.
The attempt identifies taste, customer-revenue, and belief coefficients under a specified signal and choice model, with overidentifying restrictions and an exact invariance-breaking equivalence. Defending those exclusions, handling competitive hiring capacity, and evaluating earnings and firm output remain necessary for the broader discrimination-mechanism target.

\begin{thebibliography}{9}
\bibitem{source} M. Ahmadi, G.-J. Clochard, J. Lachman and J. A. List (2025), \emph{Toward an Understanding of Discrimination When Multiple Channels Exist}. Becker Friedman Institute working-paper abstract, 29 January 2025; primary abstract checked. \url{https://bfi.uchicago.edu/working-papers/toward-an-understanding-of-discrimination-when-multiple-channels-exist/}.
\end{thebibliography}

\end{document}
